We gaan een informatieblad over Linux, Debian en KDE maken. Dit doen we in LibreOffice Writer. De titel van het document wordt: \textbf{Kennismaken met Linux KDE}

Het document moet de volgende onderdelen bevatten:
\begin{enumerate}
\item Een duidelijke titel.
\item Een korte introductie van ongeveer 100 woorden over Linux en leg uit wat het verschil is tussen Linux en Debian.
\item Een uitleg van KDE Plasma in je eigen woorden.
\item Een uitleg van LibreOffice.
\item Een afbeelding die past bij Linux of KDE. Gebruik een afbeelding die je mag gebruiken en vermeld de bron of maak zelf een screenshot (kijk eens bij de GIMP).
\item Een genummerde lijst met vijf voordelen van werken met een grafische interface.
\item Een afsluitende alinea waarin je beschrijft wat je hebt geleerd.
\end{enumerate}

Opmaakeisen
\begin{itemize}
\item Gebruik minimaal twee verschillende kopstijlen.
\item Maak de titel groter dan de gewone tekst.
\item Gebruik vetgedrukte tekst waar dat functioneel is.
\item Voeg paginanummers toe als het document meerdere pagina's beslaat.
\item Sla het document op als Kennismaken-Linux.odt in de map Documenten.
\item Exporteer het document ook naar PDF.
\end{itemize}

